\magnification=\magstep1
\input notesmac.tex
\nopagenumbers

\centerline{\titlefont Tech and Tools, Lab 2}
\bigskip
\line{Dr.~WANG Jian-Sheng \hfill for 14/16 Sep 99}
\bigskip
%----------------------------------------------------------------------- 
\problem Mathematica is good at calculus. Compute the following 
(partial) derivatives with Mathematica.
$$\halign{\h$#$\hfil&\h$#$\hfil&\h$#$\hfil\cr
{d\; \over dx} \sin x, &  {d\;\over dx} x^x & {\partial\; \over \partial y} \left(x^3-3xy^2\right), \cr
{d^2\; \over dx^2} \cos^{-1} x, & \left({\partial^2\; \over \partial x^2} + 
{\partial^2 \;\over \partial y^2}\right) \left(x+ i y\right)^3,\;\; i^2 = -1,& {d\; \over dx} f(x^2) .\cr}$$

\problem Differentiation can be done for any combination of algebraic
expressions.  But indefinite integration is much harder.  Not every
expression can be integrated.
\item{} (a) Try the following
integration by Mathematica. 
$$\halign{\h$#$\hfil&\h$#$\hfil&\h$#$\hfil\cr
\int x\; dx, & \int \sin x \;dx, & \int (1 + e^x)^{-1} \; dx. \cr}$$
(b) And some more complex (type in very very carefully):
$$\int 
{{x\,\bigl\{ 2\,x\,e^{3\,{x^2}}
        \bigl[ x - \left( 1 + x + 2\,{x^2} + 2\,{x^3} \right) \,
           \log (1 + x) \bigr]  + 
       {\bigl[ x^2 e^{2\,{x^2}} - {\log^2 (1 + x)} \bigr]^2}
        \bigr\} }\over 
   {\left( 1 + x \right) \,{{\left\{ -x^2 {e^{2\,{x^2}}} + 
          {\log^2 (1 + x)} \right\} }^2}}}
\;dx.$$
(c) Try the ratio of polynomial function. 
$$\int { 6x^5 + 6x^4-8x^3-18x^2+8x+8 \over
         x^6 - 5x^4 - 8x^3 - 2x^2 + 2x + 1 }\; dx.$$
(d) Try
$$\int { x \over 1 + e^x}\; dx$$

\problem Find the five (numerical) roots of the equation
$$ x^5 + 5 x^4 + 4 x^3 + 3 x^2 + 2 x + 1 = 0.$$
Solve the equation again with a higher precision, for
example, 30 digits of accuracy.

\problem Solve the set of second order equations
$$\eqalign{  y^2 + x^2 - y - 3x &= 0,\cr
             y^2 - 6xy - x^2 + 11y + 7x - 12 &= 0.\cr}$$
Verify that the solutions are correct by substituting the solution back.

\problem  Define a Mathematica function $f(x) = 1/(x^2+1)$, and
make a plot of it.   Define the integer factorial function $n!$
in Mathematica.  Define the greatest common divisor function
$\gcd(a,b)$ of two integers; compare with build-in function
{\tt GCD[ ...]}.  Generalize your gcd definition so that it works
for 1, 2, 3 or any number of arguments. 

\item{} Note: greatest common divisor integer $c$, of a set of integers $a_1$, 
$a_2$, $a_3$, $\ldots$, is the largest number such that $c$ divides 
$a_i$ for all $i$ exactly.  In other words, there exists an integer
$d_i$ such that $a_i = c d_i$.  Hint: recursive definition would be
best. 

\bye

